\chapter{TESTING \& EVALUATION}

\section{Testing and Evaluation Overview}

Testing and evaluation are conducted to verify that the Personalized Learning Path Platform satisfies its functional, performance, security, and AI-related requirements.

The testing process focuses on the main system functions and evaluates whether the implemented components operate correctly and support the intended learning processes.

The main testing areas include:

\begin{itemize}
	\item Functional testing.
	\item Performance testing.
	\item AI evaluation.
	\item Security testing.
	\item User and acceptance evaluation.
\end{itemize}

\section{Functional Testing}

Functional testing verifies the main business flows of the system.

The following functions should be tested:

\begin{itemize}
	\item User authentication and authorization.
	\item Competency diagnostic assessment.
	\item Learning profile generation.
	\item Personalized learning path generation.
	\item Adaptive learning recommendations.
	\item AI Tutor interaction.
	\item RAG-based question answering.
	\item Learning analytics.
	\item Examination performance prediction.
	\item Teacher monitoring.
	\item Parent monitoring.
\end{itemize}

Each function should be tested using representative test cases to verify the expected system behavior.

\subsection{Authentication and Authorization Testing}

Authentication testing verifies whether users can securely log in to the system.

Authorization testing verifies whether each user role can access only the functions permitted for that role.

The main roles include:

\begin{itemize}
	\item Student.
	\item Teacher.
	\item Parent.
	\item Administrator.
\end{itemize}

\subsection{Diagnostic Assessment Testing}

Diagnostic assessment testing verifies whether the system can:

\begin{itemize}
	\item Present assessment questions correctly.
	\item Record student answers.
	\item Calculate assessment results.
	\item Store assessment data.
	\item Use assessment results to support learning profile generation.
\end{itemize}

\subsection{Personalized Learning Testing}

Personalized learning testing verifies whether the system can generate learning recommendations and learning paths based on available student information.

The evaluation should consider factors such as:

\begin{itemize}
	\item Competency level.
	\item Assessment results.
	\item Learning progress.
	\item Learning activities.
	\item Learning needs.
\end{itemize}

\section{Performance Testing}

Performance testing evaluates whether the system can satisfy the proposed non-functional performance requirements.

The proposal specifies an average response time of less than 3 seconds for standard operations and support for at least 1,000 concurrent active users.

Therefore, performance testing should measure:

\begin{itemize}
	\item Average API response time.
	\item Concurrent-user capacity.
	\item Database response performance.
	\item AI service latency.
	\item System behavior under increased workload.
\end{itemize}

The performance test results should be recorded and compared with the defined requirements.

If the actual system has not yet been tested under the required workload, no performance result should be reported as an achieved result.

\section{AI Evaluation}

AI components require dedicated evaluation because their outputs cannot be evaluated only through conventional functional testing.

\subsection{Recommendation System Evaluation}

The Recommendation Engine should be evaluated based on the relevance of recommended learning content to the student's learning profile, competency level, and learning progress.

\subsection{Knowledge Tracing Evaluation}

Knowledge Tracing should be evaluated based on its ability to consistently estimate changes in students' knowledge or mastery using assessment and learning activity data.

\subsection{Performance Prediction Evaluation}

The Student Performance Prediction component should be evaluated by comparing predicted examination performance with appropriate evaluation data.

The final evaluation metrics should depend on the prediction model and dataset actually used by the development team.

\subsection{AI Tutor Evaluation}

The AI Tutor should be evaluated based on:

\begin{itemize}
	\item Answer relevance.
	\item Answer quality.
	\item Academic correctness.
	\item Response consistency.
	\item Unsupported or hallucinated information.
\end{itemize}

\subsection{RAG Evaluation}

The RAG-based question-answering system should be evaluated based on:

\begin{itemize}
	\item Retrieval relevance.
	\item Retrieved information quality.
	\item Relevance of generated answers.
	\item Grounding of answers in the retrieved knowledge.
	\item Unsupported-answer rate.
\end{itemize}

\section{Security Testing}

Security testing verifies whether the system protects user accounts, learning data, and APIs.

The main security testing areas include:

\begin{itemize}
	\item Authentication.
	\item Role-Based Access Control.
	\item API access control.
	\item Input validation.
	\item Data protection.
	\item Secure communication.
	\item Backup and recovery procedures.
\end{itemize}

Security testing should verify that unauthorized users cannot access protected functions or data.

\section{Learning Analytics Evaluation}

The Learning Analytics component should be evaluated to verify whether the system correctly collects and presents learning information.

The evaluation should consider:

\begin{itemize}
	\item Accuracy of recorded learning activities.
	\item Accuracy of assessment results.
	\item Correctness of progress information.
	\item Correctness of displayed performance information.
	\item Accessibility of dashboards for different user roles.
\end{itemize}

\section{Usability Evaluation}

Usability evaluation focuses on whether students, teachers, and parents can use the system effectively.

The evaluation may consider:

\begin{itemize}
	\item Ease of navigation.
	\item Clarity of information.
	\item Ease of completing learning activities.
	\item Ease of monitoring learning progress.
	\item Overall user experience.
\end{itemize}

\section{Acceptance Criteria}

The system should be considered ready for final evaluation when the major project requirements have been implemented and evaluated.

The main acceptance criteria include:

\begin{itemize}
	\item Core business flows are implemented.
	\item Major functional test cases pass.
	\item Performance requirements are measured.
	\item AI components are evaluated using appropriate metrics.
	\item Security requirements are tested.
	\item Learning analytics functions operate correctly.
	\item Required project deliverables are demonstrated.
\end{itemize}

\section{Evaluation Status}

At the current stage, the testing and evaluation chapter defines the testing strategy and evaluation criteria for the proposed system.

Actual test results, performance measurements, AI evaluation metrics, screenshots, and acceptance results should be added after the corresponding system components have been implemented and tested by the development team.